\documentclass[12pt]{article}
\usepackage[T1]{fontenc}
\usepackage[utf8]{inputenc}
\usepackage{lmodern}
\usepackage{textcomp}
\usepackage{enumitem}
\usepackage{geometry}
\geometry{margin=1in}
\begin{document}
\begin{center}
Answer Key - Health Sciences - Undergraduate Level Assessment 20 \
\small (Answers are ordered exactly as in the question paper.)
\end{center}

\section*{Multiple Choice}
\begin{enumerate}[label=\arabic*.]
\item \textbf{Answer:} C
\\ \textit{Options:} A) Father; B) Mother; C) Sister; D) Son
\\ \textit{Note:} An individual is most likely to share a common HLA haplotype with their sister because siblings inherit one haplotype from each parent, leading to a greater probability of sharing identical haplotypes compared to more distant relatives.
\item \textbf{Answer:} B
\\ \textit{Options:} A) Pope John II's; B) Pope Paul VI's; C) Pope John Paul I's; D) Pope John Paul II's
\\ \textit{Note:} The correct answer is "Pope Paul VI's" because Humanae Vitae is the encyclical issued by Pope Paul VI in 1968 that explicitly condemns artificial birth control, articulating the Catholic Church's teachings on human sexuality and procreation.
\item \textbf{Answer:} D
\\ \textit{Options:} A) Technohouse; B) Compuhouse; C) Futurehouse; D) Smarthouse
\\ \textit{Note:} The term "Smarthouse" accurately describes homes equipped with advanced technology that enables centralized control of various systems and appliances, exemplifying the integration of smart technology in modern living environments.
\item \textbf{Answer:} A
\\ \textit{Options:} A) provide a lot of reinforcement and a small amount of punishment; B) provide more punishment than reinforcement; C) provide an equal balance of punishment and reinforcement; D) are extremely attractive and provide a lot of reinforcement
\\ \textit{Note:} Byrne's law of attraction posits that individuals are drawn to others who offer positive reinforcement, such as support and approval, while minimizing negative experiences or punishment, creating a favorable and balanced social dynamic.
\item \textbf{Answer:} B
\\ \textit{Options:} A) Sensory memory; B) Episodic memory; C) Semantic memory; D) Procedural memory
\\ \textit{Note:} Episodic memory, which involves the recall of specific events and experiences, shows the most significant decline with advanced age due to the natural deterioration of cognitive processes associated with retrieving personal memories.
\end{enumerate}

\section*{True / False}
\begin{enumerate}[label=\arabic*.]
\setcounter{enumi}{5}
\item \textbf{Answer:} False
\\ \textit{Options:} A) True; B) False
\\ \textit{Note:} Glucose-6-phosphate dehydrogenase (G$_{6}$PD) deficiency is inherited in an X-linked recessive pattern, not an autosomal dominant pattern, primarily affecting males who inherit the gene mutation from their carrier mothers.
\item \textbf{Answer:} False
\\ \textit{Options:} A) True; B) False
\\ \textit{Note:} Maggie Kuhn was a prominent activist known for founding the Gray Panthers, but she did not live to be 115 years old, as the oldest verified person, Jeanne Calment, reached that age.
\item \textbf{Answer:} True
\\ \textit{Options:} A) True; B) False
\\ \textit{Note:} The statement is true because during the pharyngeal phase of swallowing, the muscles around the larynx contract to close the laryngeal entrance, effectively preventing food from entering the airway and ensuring it is directed into the esophagus.
\item \textbf{Answer:} True
\\ \textit{Options:} A) True; B) False
\\ \textit{Note:} Open-ended questions on a survey often lead to diverse and varied responses that can be difficult to categorize and analyze, complicating the interpretation of the answers provided by respondents.
\item \textbf{Answer:} False
\\ \textit{Options:} A) True; B) False
\\ \textit{Note:} A person with a rating of 4 on the Kinsey scale is actually classified as predominantly homosexual with incidental heterosexual experience, not mostly heterosexual.
\end{enumerate}

\section*{Long Form Response}
\begin{enumerate}[label=\arabic*.]
\setcounter{enumi}{10}
\item \textbf{Answer:} The MMR vaccine, which protects against measles, mumps, and rubella, plays a crucial role in controlling measles outbreaks globally through its high effectiveness in inducing immunity. Vaccination coverage is essential for achieving herd immunity; however, challenges such as vaccine hesitancy, misinformation, and access disparities hinder widespread uptake. Despite these obstacles, the MMR vaccine has significantly reduced the incidence of measles and related complications, proving vital for public health efforts. The resurgence of measles in some regions underscores the need for consistent vaccination campaigns and education to maintain high coverage levels and protect vulnerable populations. Ultimately, the MMR vaccine remains a cornerstone of measles prevention strategies worldwide.
\\ \textit{Note:} correct because it effectively highlights the MMR vaccine's high effectiveness in preventing measles, identifies the key challenges to achieving adequate vaccination coverage, and emphasizes its critical role in reducing measles incidence and improving public health outcomes globally.
\item \textbf{Answer:} Cocaine activates the sympathetic nervous system primarily by inhibiting the reuptake of neurotransmitters such as dopamine, norepinephrine, and serotonin in the brain. This leads to an increase in the levels of these neurotransmitters, resulting in heightened arousal and stimulation of the "fight or flight" response. However, while this activation can initially enhance energy and alertness, it also causes physiological responses such as increased heart rate and blood pressure, which can detract from sexual performance. The increased sympathetic activity can lead to anxiety and reduced blood flow to the genital area, ultimately impairing sexual function. The implications for health and behavior are significant, as chronic cocaine use can lead to sexual dysfunction, relationship issues, and broader health complications, highlighting the need for awareness and intervention in substance use disorders.
\\ \textit{Note:} correct because it accurately describes how cocaine inhibits neurotransmitter reuptake, leading to sympathetic nervous system activation and associated physiological responses that can hinder sexual performance while initially increasing arousal and energy levels.
\end{enumerate}

\vfill
\noindent\textit{End of Answer Key}
\end{document}